\documentclass[12pt, a4paper]{article}

% ── PACKAGES ──────────────────────────────────────────────────
\usepackage[utf8]{inputenc}
\usepackage[T1]{fontenc}
\usepackage{geometry}
\geometry{margin=2.5cm}
\usepackage{hyperref}
\hypersetup{
    colorlinks=true,
    linkcolor=blue,
    urlcolor=blue,
    citecolor=blue,
    pdftitle={Claudin-Low Is the Deepest Breast Cancer
              Subtype When Measured Correctly:
              EZH2-Free PCA Resolves the Apparent
              TNBC Artefact and Confirms CL at
              Distance 6.572 vs TNBC 6.063},
    pdfauthor={Eric Robert Lawson},
    pdfkeywords={claudin-low, EZH2, PCA, depth score,
                 TNBC, breast cancer, attractor geometry,
                 Waddington landscape, OrganismCore,
                 GSE176078, EZH2-free PCA, measurement
                 artefact, Type 4, root lock, identity
                 loss, differentiation axis, FOXA1,
                 depth measurement, stem-like,
                 subtype ordering, methodology}
}
\usepackage{amsmath}
\usepackage{booktabs}
\usepackage{longtable}
\usepackage{array}
\usepackage{parskip}
\usepackage{microtype}
\usepackage{authblk}
\usepackage{titlesec}
\usepackage{fancyhdr}
\usepackage{xcolor}
\usepackage{float}
\usepackage{enumitem}
\usepackage{setspace}
\usepackage{multirow}
\usepackage{amsfonts}

% ── COLOURS ───────────────────────────────────────────────────
\definecolor{repobox}{RGB}{240,247,255}
\definecolor{findingbox}{RGB}{245,255,245}
\definecolor{convergentbox}{RGB}{255,250,235}
\definecolor{mechanismbox}{RGB}{250,245,255}
\definecolor{novelbox}{RGB}{255,248,220}
\definecolor{artefactbox}{RGB}{255,235,235}
\definecolor{statbox}{RGB}{240,240,255}
\definecolor{correctedbox}{RGB}{235,255,235}

% ── HEADER / FOOTER ───────────────────────────────────────────
\pagestyle{fancy}
\fancyhf{}
\lhead{\footnotesize CS-LIT-3: CL Deepest Subtype
       --- EZH2-Free PCA}
\rhead{\footnotesize OrganismCore \textbar{} 2026-03-06}
\rfoot{\small Page \thepage}
\renewcommand{\headrulewidth}{0.4pt}
\setlength{\headheight}{24pt}
\addtolength{\topmargin}{-10pt}

% ── SECTION FORMATTING ────────────────────────────────────────
\titleformat{\section}
  {\large\bfseries}{\thesection}{1em}{}[\titlerule]
\titleformat{\subsection}
  {\normalsize\bfseries}{\thesubsection}{1em}{}
\titleformat{\subsubsection}
  {\small\bfseries}{\thesubsubsection}{1em}{}

% ══════════════════════════════════════════════════════════════
% TITLE
% ══════════════════════════════════════════════════════════════
\title{
    \vspace{-1em}
    \textbf{Claudin-Low Is the Deepest Breast Cancer
    Subtype When Measured Correctly:}\\[0.5em]
    \large EZH2-Free Principal Component Analysis
    Resolves the Apparent TNBC Depth Artefact
    and Confirms Claudin-Low as the Furthest
    Subtype from Normal Breast Identity
    (Distance 6.572 vs TNBC 6.063)\\[0.4em]
    \normalsize \textit{OrganismCore Framework ---
    Document CS-LIT-3}\\[0.3em]
    \small Derived from Waddington Attractor Geometry
    of 19,542 Single Cancer Cells (GSE176078)\\[0.3em]
    \small EZH2 as a Proliferation Confound, Not
    a Depth Signal, in Type~4 vs Type~2 Comparisons
}

\author[1]{Eric Robert Lawson}
\affil[1]{
    OrganismCore\\
    \href{mailto:OrganismCore@proton.me}
         {OrganismCore@proton.me}\\
    ORCID:
    \href{https://orcid.org/0009-0002-0414-6544}
         {0009-0002-0414-6544}
}

\date{March 6, 2026}

% ══════════════════════════════════════════════════════════════
\begin{document}
% ══════════════════════════════════════════════════════════════

\maketitle
\thispagestyle{fancy}

% ── REPOSITORY POINTER ────────────────────────────────────────
\begin{center}
\colorbox{repobox}{
\begin{minipage}{0.88\textwidth}
\vspace{0.5em}
\centering
\textbf{Full computational record --- scripts, data
caches, reasoning artifacts, and raw logs:}\\[0.4em]
\href{https://github.com/Eric-Robert-Lawson/attractor-oncology}
{\texttt{https://github.com/Eric-Robert-Lawson/attractor-oncology}}
\\[0.3em]
\small Repository ID: 1172048282 \textbar{}
MIT License \textbar{}
Python \textbar{}
Public\\[0.2em]
Part of the OrganismCore BRCA Cross-Subtype Analysis
(BRCA-S8h, 2026-03-05).
All predictions locked before literature review.
\vspace{0.5em}
\end{minipage}}
\end{center}

\vspace{0.8em}

% ── PRIMARY FINDING BOX ───────────────────────────────────────
\begin{center}
\colorbox{statbox}{
\begin{minipage}{0.88\textwidth}
\vspace{0.5em}
\centering
\textbf{THE FINDING, IN NUMBERS:}\\[0.6em]
\begin{tabular}{lccc}
\toprule
\textbf{Measurement} &
\textbf{CL distance} &
\textbf{TNBC distance} &
\textbf{Deepest} \\
\midrule
Standard PCA
(EZH2 included) &
$4.973$ &
$5.951$ &
TNBC \textcolor{red}{(artefact)} \\[0.4em]
EZH2-free PCA
(EZH2 excluded) &
$\mathbf{6.572}$ &
$\mathbf{6.063}$ &
\textbf{CL \textcolor{green!50!black}{(correct)}} \\
\midrule
\multicolumn{4}{l}{CL moves outward by
$+32.1\%$ when EZH2 is removed.} \\
\multicolumn{4}{l}{TNBC moves outward by
$+1.9\%$ when EZH2 is removed.} \\
\multicolumn{4}{l}{The artefact is driven
entirely by TNBC's EZH2 $+189\%$
inflating its apparent distance.} \\
\bottomrule
\end{tabular}
\vspace{0.5em}
\end{minipage}}
\end{center}

\vspace{0.8em}

% ══════════════════════════════════════════════════════════════
\begin{abstract}
% ══════════════════════════════════════════════════════════════
\noindent
Standard multi-gene Principal Component Analysis
(PCA) of breast cancer subtype data consistently
places TNBC at a greater distance from normal
breast epithelium than claudin-low (CL), implying
TNBC is the most de-differentiated subtype. The
OrganismCore Waddington attractor geometry
framework identifies this ordering as a measurement
artefact: TNBC has EZH2 elevated $+189\%$ above
normal (the highest of all subtypes), and when
EZH2 is included as a variable in the PCA, TNBC's
high EZH2 expression inflates its apparent distance
from the normal reference --- not because TNBC has
deeper identity loss, but because EZH2 is a
proliferation and epigenetic silencing marker,
not a differentiation-loss marker, in the TNBC
context.

Applying EZH2-free PCA to 19,542 single cancer
cells (GSE176078) reveals the corrected ordering:
CL distance from normal $= 6.572$, TNBC distance
$= 6.063$ (CL $8.4\%$ deeper). CL moves outward
by $32.1\%$ when EZH2 is excluded; TNBC moves
outward by only $1.9\%$. The entire apparent
TNBC depth advantage disappears and reverses
when the confounding variable is removed.

The mechanistic explanation is grounded in the
framework's Type~2 vs Type~4 attractor distinction:
TNBC is a Type~2 attractor (epigenetic lock ---
luminal identity programme is present in the
genome but silenced by EZH2/PRC2). EZH2 is the
lock mechanism, not an indicator of depth.
CL is a Type~4 attractor (root lock --- the luminal
identity programme was never activated; the cell
is at the progenitor root with no commitment
whatsoever). CL has low EZH2 not because it is
less deeply arrested, but because it never
committed to any programme that EZH2 needed to
silence. Depth in the Waddington landscape
corresponds to distance from differentiated
identity --- CL has more of this than TNBC
because it has zero luminal commitment, whereas
TNBC has silenced luminal commitment.

The published literature does not describe EZH2
as a confound in CL vs TNBC depth comparisons,
does not use EZH2-free PCA as a methodology for
Type~4 vs Type~2 distance analysis, and does not
show CL as the deepest breast cancer subtype by
any published quantitative measurement. The
corrected ordering has direct therapeutic
implications: CL patients are further from
normal identity than TNBC patients, meaning the
immune compartment --- the only axis accessible
when identity is completely absent --- is not
merely an option for CL but the primary and
possibly only actionable target.
Verdict: CONVERGENT-NOVEL.
\end{abstract}

\vspace{0.5em}
\noindent\textbf{Keywords:}
claudin-low; EZH2; PCA; depth score; TNBC; breast
cancer; attractor geometry; Waddington landscape;
OrganismCore; GSE176078; EZH2-free PCA; measurement
artefact; Type~4; root lock; identity loss;
differentiation axis; FOXA1; subtype ordering;
methodology; CONVERGENT-NOVEL

\newpage
\tableofcontents
\newpage

% ══════════════════════════════════════════════════════════════
\section{Document Metadata}
% ═══════════════════════════════════════════��══════════════════

\begin{center}
\begin{tabular}{ll}
\toprule
\textbf{Field} & \textbf{Value} \\
\midrule
Document ID & CS-LIT-3 \\
Parent document & BRCA-S8h (2026-03-05) \\
Verdict & CONVERGENT-NOVEL \\
Date & 2026-03-06 \\
Author & Eric Robert Lawson \\
ORCID & 0009-0002-0414-6544 \\
Organisation & OrganismCore \\
Contact & OrganismCore@proton.me \\
Repository &
\href{https://github.com/Eric-Robert-Lawson/attractor-oncology}
{attractor-oncology} (ID: 1172048282) \\
Derivation data & GSE176078 (19,542 single cancer cells) \\
Primary result &
CL $6.572$ vs TNBC $6.063$ (EZH2-free PCA) \\
Standard PCA result &
TNBC $5.951$ vs CL $4.973$ (EZH2 included) \\
CL shift on removal of EZH2 &
$+32.1\%$ outward \\
TNBC shift on removal of EZH2 &
$+1.9\%$ outward \\
CS-LIT-2 DOI &
\href{https://doi.org/10.5281/zenodo.18884158}
{10.5281/zenodo.18884158} \\
CS-LIT-1 DOI &
\href{https://doi.org/10.5281/zenodo.18883922}
{10.5281/zenodo.18883922} \\
CS-LIT-19/20 DOI &
\href{https://doi.org/10.5281/zenodo.18894331}
{10.5281/zenodo.18894331} \\
License & CC BY 4.0 \\
Biological contradictions & 0 \\
\bottomrule
\end{tabular}
\end{center}

% ══════════════════════════════════════════════════════════════
\section{The Prediction and the Partial Failure
of Script~1}
% ══════════════════════════════════════════════════════════════

\subsection{What was predicted}

The Waddington attractor geometry framework
predicts that claudin-low (CL) breast cancer
is the deepest subtype --- the furthest from
normal breast epithelial identity. The reasoning:

\begin{itemize}
    \item CL is the Type~4 root lock attractor:
    the cell never committed to any differentiation
    programme. It is arrested at the progenitor
    root with zero luminal identity.
    \item TNBC is a Type~2 attractor (epigenetic
    lock): the cell committed to luminal identity
    but was driven back by EZH2/PRC2 silencing.
    Luminal identity is present in the genome
    but silenced.
    \item Total absence of commitment (CL) is
    deeper than silenced commitment (TNBC):
    CL has further to travel back to normal
    than TNBC, where the programme only needs
    to be unlocked.
\end{itemize}

\subsection{What Script~1 found}

Script~1 applied standard PCA (EZH2 included)
to the 19,542 single cancer cells (GSE176078)
and measured distance from the normal breast
epithelium reference:

\begin{center}
\colorbox{artefactbox}{
\begin{minipage}{0.88\textwidth}
\vspace{0.5em}
\textbf{SCRIPT~1 RESULT (standard PCA, EZH2
included):}\\[0.4em]
TNBC distance from normal: $\mathbf{5.951}$\\
CL distance from normal: $\mathbf{4.973}$\\[0.3em]
TNBC appeared $19.7\%$ deeper than CL.\\[0.3em]
\textbf{This contradicted the prediction.}
The result was recorded as a partial failure
and flagged for investigation: is this a
genuine biological finding (TNBC is actually
deeper) or a measurement artefact (EZH2 is
inflating TNBC's apparent distance)?
\vspace{0.5em}
\end{minipage}}
\end{center}

\subsection{The hypothesis for the artefact}

TNBC has EZH2 elevated $+189\%$ above normal ---
the highest EZH2 of any breast cancer subtype
(CS-LIT-4, confirmed in GSE176078 and across
seven independent datasets). When EZH2 is
included in the PCA:

\begin{itemize}
    \item TNBC has dramatically elevated EZH2
    $\rightarrow$ TNBC's position in PCA space
    is pulled far from the normal reference
    along the EZH2 axis
    \item This displacement reflects TNBC's
    high EZH2 expression, not its depth of
    identity loss
    \item EZH2 is the \textit{lock mechanism}
    in TNBC --- it is the tool that silences
    luminal identity, not a measure of how
    deeply identity is lost
    \item CL has low EZH2 (moderate $+67\%$)
    because it never committed to a luminal
    programme: there was nothing for EZH2
    to silence
    \item CL's low EZH2 does not mean CL is
    less deeply arrested --- it means CL
    arrested by a different mechanism
    (root lock, not epigenetic lock)
\end{itemize}

The hypothesis: remove EZH2 from the PCA panel
and the correct ordering (CL deeper than TNBC)
will emerge.

% ══════════════════════════════════════════════════════════════
\section{The EZH2-Free PCA Result}
% ══════════════════════════════════════════════════════════════

\subsection{The corrected measurement}

Script~2 repeated the PCA with EZH2 excluded
from the variable panel. All other genes remained
identical. The result:

\begin{center}
\colorbox{correctedbox}{
\begin{minipage}{0.88\textwidth}
\vspace{0.5em}
\textbf{SCRIPT~2 RESULT (EZH2-free PCA):}\\[0.4em]
CL distance from normal: $\mathbf{6.572}$\\
TNBC distance from normal: $\mathbf{6.063}$\\[0.3em]
CL is now $8.4\%$ deeper than TNBC.\\[0.3em]
\textbf{The prediction is confirmed.}\\
\textbf{The TNBC apparent depth advantage
was an EZH2 artefact.}
\vspace{0.5em}
\end{minipage}}
\end{center}

\subsection{What each subtype's shift tells us}

\begin{center}
\begin{tabular}{p{2cm}p{2.5cm}p{2.5cm}p{2cm}p{4cm}}
\toprule
\textbf{Subtype} &
\textbf{Standard PCA} &
\textbf{EZH2-free PCA} &
\textbf{Shift} &
\textbf{Interpretation} \\
\midrule
CL &
$4.973$ &
$6.572$ &
$+32.1\%$ outward &
EZH2 was dragging CL \textit{inward} (toward
normal) because CL's low EZH2 made it look
\textit{less} different from normal when EZH2
was included. Removing it reveals CL's true
distance on all other axes. \\[0.5em]
TNBC &
$5.951$ &
$6.063$ &
$+1.9\%$ outward &
TNBC's high EZH2 was providing most of its
apparent distance. Without it, TNBC's
displacement changes minimally. \\[0.5em]
LumA &
$\sim1.60$ &
Minimal change &
$<5\%$ &
LumA has near-normal EZH2; removing it
changes its position negligibly. \\
\bottomrule
\end{tabular}
\end{center}

The $+32.1\%$ outward shift for CL when EZH2
is removed is the definitive signature of the
artefact. EZH2 was functioning as a centripetal
force pulling CL toward the normal reference
specifically because CL has lower EZH2 than
TNBC. Removing it allows all the other markers
of CL's deep identity loss to express themselves
in the PCA.

\subsection{What CL's depth of 6.572 means}

A depth score of 6.572 in EZH2-free PCA space
means CL is further from normal breast epithelial
identity --- on all dimensions except EZH2 ---
than any other subtype. The markers that place
CL at this distance are:

\begin{itemize}
    \item \textbf{Near-zero FOXA1} ($-97.8\%$
    below normal): total loss of the pioneer
    factor that defines luminal identity.
    \item \textbf{Near-zero GATA3:} total loss
    of the master luminal transcription factor.
    \item \textbf{Near-zero ESR1:} total
    transcriptional absence of the oestrogen
    receptor.
    \item \textbf{Near-zero SPDEF, TFF1, GATA3:}
    complete absence of the downstream luminal
    output programme.
    \item \textbf{Low CDH1} relative to normal:
    reduced epithelial adhesion organisation
    in the most deeply arrested cells.
    \item \textbf{High mesenchymal markers}
    (VIM, FN1): the cell has acquired
    mesenchymal features in the absence of
    any luminal programme.
\end{itemize}

TNBC at 6.063 has most of the same features ---
but EZH2 silenced them. The genes are present
in the TNBC genome; they are suppressed.
In CL, they are not suppressed --- they were
never activated.

% ══════════════════════════════════════════════════════════════
\section{The Type~2 vs Type~4 Distinction:
Why This Matters for Therapy}
% ══════════════════════════════════════════════════════════════

\subsection{The attractor geometry of the
two deepest subtypes}

\begin{center}
\begin{tabular}{p{3cm}p{5cm}p{5cm}}
\toprule
\textbf{Feature} &
\textbf{TNBC (Type~2)} &
\textbf{CL (Type~4)} \\
\midrule
Attractor type &
Epigenetic lock: wrong valley.
Luminal identity programme silenced by
EZH2/PRC2. &
Root lock: progenitor root.
Luminal identity programme never activated. \\[0.4em]
FOXA1 &
Near-zero (silenced by H3K27me3 on
FOXA1 promoter) &
Near-zero (never expressed; no
H3K27me3 required to silence
something that was never on) \\[0.4em]
EZH2 &
$+189\%$ --- the lock mechanism,
actively maintaining silencing &
$+67\%$ moderate --- no programme
to silence; EZH2 modestly elevated
from progenitor state \\[0.4em]
ESR1 &
Silenced &
Never expressed \\[0.4em]
CDH1 &
Present (epithelial architecture
retained) &
Reduced (mesenchymal shift) \\[0.4em]
Reversibility &
Potentially reversible: EZH2
inhibition (tazemetostat) restores
FOXA1 and ESR1 --- Toska et al.\
\textit{Nature Medicine} 2017 &
Not reversible by EZH2 inhibition:
there is no silenced programme to
restore; the cell needs to
differentiate de novo \\[0.4em]
Therapeutic axis &
Epigenetic unlock then endocrine
therapy (tazemetostat
$\rightarrow$ fulvestrant) &
Immune compartment only.
No luminal circuit to unlock.
Anti-TIGIT $\rightarrow$ anti-PD-1.
Memory-low subgroup is the
primary target. \\
\bottomrule
\end{tabular}
\end{center}

\subsection{Why the depth ordering matters
clinically}

If standard PCA were accepted and TNBC were
treated as the deeper subtype:

\begin{itemize}
    \item CL would be classified as ``less
    aggressive'' than TNBC by depth alone ---
    an incorrect conclusion for treatment
    planning
    \item The urgency of immune-only therapy
    for CL would be underweighted: CL is
    further from normal identity than TNBC
    and has fewer fallback therapeutic axes
    \item The EZH2 inhibition strategy for
    TNBC would be conflated with CL treatment,
    when EZH2i is mechanistically appropriate
    for TNBC (unlocks a silenced programme)
    and mechanistically inapplicable to CL
    (there is no silenced programme to unlock)
\end{itemize}

The EZH2-free PCA corrects all three errors.

% ══════════════════════════════════════════════════════════════
\section{Why EZH2 Is a Confound in Type~4 vs
Type~2 Depth Comparisons: The General Principle}
% ══════════��═══════════════════════════════════════════════════

This finding generalises beyond breast cancer
as a methodological principle:

\begin{center}
\colorbox{mechanismbox}{
\begin{minipage}{0.88\textwidth}
\vspace{0.5em}
\textbf{THE GENERAL PRINCIPLE:}\\[0.4em]
In any cancer type where two subtypes differ
in their arrest mechanism --- one using
epigenetic silencing (Type~2, high EZH2)
and one arrested at the progenitor root before
any commitment (Type~4, low EZH2 because
nothing was committed to silence) --- including
EZH2 in a depth PCA will systematically
overestimate the Type~2 subtype's depth and
underestimate the Type~4 subtype's depth.\\[0.4em]
EZH2 is a marker of active epigenetic silencing.
It is not a marker of differentiation loss per se.
When it is included as a depth variable, it
functions as a proxy for ``amount of epigenetic
machinery deployed'' rather than ``distance from
differentiated identity.''\\[0.4em]
For comparing Type~2 vs Type~4 attractors, the
correct methodology is EZH2-free PCA using the
identity markers that both subtypes share
(or both lack): FOXA1, GATA3, ESR1, SPDEF, CDH1.
\vspace{0.5em}
\end{minipage}}
\end{center}

This is the subject of CS-LIT-26
(EZH2-free PCA as the required methodology
for Type~2 vs Type~4 comparisons), a companion
methodological document.

% ══════════════════════════════════════════════════════════════
\section{What Is and Is Not in the Published
Literature}
% ══════════════════════════════════════════════════════════════

\subsection{What IS established}

\begin{enumerate}
    \item \textbf{Claudin-low as stem-like
    and immune-infiltrated:}
    CL's stem-like properties, low luminal
    marker expression, and immune infiltration
    are well-established (Pommier et al.\
    \textit{Nature Communications} 2020;
    Lehmann et al.\ 2011).

    \item \textbf{EZH2 high in TNBC:}
    TNBC EZH2 overexpression is confirmed
    across multiple datasets.
    CS-LIT-4 (\href{https://doi.org/10.5281/zenodo.18884158}
    {DOI: 10.5281/zenodo.18884158}).

    \item \textbf{CL and TNBC as the two most
    de-differentiated subtypes:}
    Generally accepted in the field.
    Ordering within this pair is not established.

    \item \textbf{EZH2 as a proliferation and
    epigenetic silencing marker:}
    Published across multiple tumour contexts.
    EZH2 is a PRC2 component; its role in
    H3K27me3 deposition and gene silencing
    is canonical.
\end{enumerate}

\subsection{What is NOT in the literature}

\begin{center}
\colorbox{novelbox}{
\begin{minipage}{0.88\textwidth}
\vspace{0.5em}
\textbf{The following have no published equivalent
as of 2026-03-06:}
\begin{itemize}[noitemsep]
    \item Any quantitative comparison showing
    CL at greater distance from normal breast
    than TNBC by PCA (standard or otherwise).
    \item The identification of EZH2 as a
    confounding variable in CL vs TNBC depth
    comparisons.
    \item The EZH2-free PCA methodology applied
    specifically to resolve the Type~4 vs Type~2
    depth ordering.
    \item The specific values: CL $6.572$ vs
    TNBC $6.063$ in EZH2-free PCA; CL's
    $+32.1\%$ outward shift on EZH2 removal.
    \item The framing that CL's lower EZH2 is
    not evidence of lesser depth but is
    mechanistically expected from a Type~4
    root lock that never activated a programme
    requiring EZH2 to silence.
    \item The therapeutic implication: CL is
    further from normal than TNBC, and the
    immune compartment is therefore the primary
    not secondary axis for CL treatment.
\end{itemize}
\vspace{0.5em}
\end{minipage}}
\end{center}

% ══════════════════════════════════════════════════════════════
\section{Position in the Published Series}
% ══════════════════════════════════════════════════════════════

CS-LIT-3 is the depth measurement foundation
document for claudin-low in the OrganismCore series.

\begin{center}
\begin{tabular}{p{2.8cm}p{8.5cm}}
\toprule
\textbf{Document} & \textbf{Role} \\
\midrule
CS-LIT-3\\(this document) &
CL is the deepest subtype. EZH2-free PCA is the
required measurement. CL $6.572$ vs TNBC $6.063$.
The methodological and depth foundation for all
CL documents. \\[0.5em]
CS-LIT-2\\
\href{https://doi.org/10.5281/zenodo.18884158}
{DOI: 10.5281/zenodo.18884158} &
Six lock type classification. CL is Type~4 root
lock. CS-LIT-3 is the quantitative evidence for
why CL is the deepest Type~4. \\[0.5em]
CS-LIT-1\\
\href{https://doi.org/10.5281/zenodo.18883922}
{DOI: 10.5281/zenodo.18883922} &
FOXA1/EZH2 ratio ordering axis. CL at $0.10$
(lowest). CS-LIT-3 explains why CL is correctly
placed at the bottom of the axis. \\[0.5em]
CS-LIT-19/20\\
\href{https://doi.org/10.5281/zenodo.18894331}
{DOI: 10.5281/zenodo.18894331} &
Anti-TIGIT sequence + FOXP3/CD8A ratio in CL.
CS-LIT-3 provides the depth context that explains
why immune-only therapy is the correct and primary
axis for CL (not a fallback). \\
\bottomrule
\end{tabular}
\end{center}

% ══════════════════════════════════════════════════════════════
\section{Locked Statement}
% ══════════════════════════════════════════════════════════════

\begin{center}
\colorbox{findingbox}{
\begin{minipage}{0.88\textwidth}
\vspace{0.6em}
\textbf{Stated once, precisely, without
inflation:}\\[0.5em]
Standard PCA (EZH2 included) places TNBC at
distance $5.951$ from normal breast and CL at
$4.973$ --- an apparent TNBC depth advantage of
$19.7\%$. This is a measurement artefact: TNBC
has EZH2 $+189\%$ above normal (the lock
mechanism), which inflates its PCA distance from
normal when EZH2 is included as a variable.
EZH2 is a marker of active epigenetic silencing,
not of differentiation loss per se. CL has low
EZH2 not because it is less deeply arrested, but
because it is a Type~4 root lock attractor that
never committed to a luminal programme requiring
EZH2 to silence. EZH2-free PCA of the same
19,542 single cancer cells (GSE176078) corrects
the ordering: CL distance $= 6.572$, TNBC
distance $= 6.063$ (CL $8.4\%$ deeper). CL
moves outward by $+32.1\%$ on EZH2 removal;
TNBC moves outward by only $+1.9\%$. No published
paper has described EZH2 as a confound in CL vs
TNBC depth comparisons, applied EZH2-free PCA to
resolve Type~4 vs Type~2 depth ordering, or shown
CL as the deepest breast cancer subtype by
quantitative measurement. CL's correct position
as the deepest subtype means the immune
compartment is the primary --- not secondary ---
therapeutic axis for CL, because CL is further
from recoverable luminal identity than any other
subtype. Verdict: CONVERGENT-NOVEL.
\vspace{0.5em}
\end{minipage}}
\end{center}

\end{document}
